% !PRINTED 18
\weberpage{18}
%<<CONT
\begingroup
\setlength{\abovedisplayskip}{5pt plus 2pt minus 1pt}
\setlength{\belowdisplayskip}{5pt plus 2pt minus 1pt}
\setlength{\abovedisplayshortskip}{3pt plus 1pt}
\setlength{\belowdisplayshortskip}{4pt plus 1pt}
und wir setzen ausserdem fest, dass $(x,0)=x$ sei, was diesen
Gleichungen nicht widerspricht. Es ist $(x,y)$ nur dann $=0$, wenn $x$
und $y$ beide gleich Null sind. Ferner bezeichnen wir zur Abkürzung
$(0,1)$ mit $i$. Dann ergeben obige Gleichungen die Folgerungen:
\[
 (x,0)(0,1)=(0,x)\quad\text{oder}\quad =ix,
\]
\[
 (x,0)+(0,y)=(x,y)=x+yi,
\]
\[
 x+yi+a+bi=(x+a)+(y+b)i,
\]
und die Umkehrung der Addition:
\[
 x+yi-(a+bi)=(x-a)+i(y-b),
\]
ferner die Multiplication:
\[
 (x+yi)(a+bi)=xa-yb+i(xb+ya),
 \qquad i^2=-1,
\]
\[
 (x+yi)(x-yi)=x^2+y^2,
\]
\[
 x+yi=(a+bi){(a-bi)(x+yi)\over a^2+b^2}
\]
oder
\[
 {x+yi\over a+bi}={ax+by+i(ay-bx)\over a^2+b^2},
\]
wodurch die Division erklärt ist, ausser wenn $a+bi=0$ ist.

Zahlen von der Form $ix$ heissen \gesperrt{rein imaginäre Zahlen} und
$i$ die imaginäre Einheit. Die Zahlen $a+bi$ heissen
\gesperrt{imaginär} oder \gesperrt{complex}. Das System der
\gesperrt{reellen} und der \gesperrt{rein imaginären Zahlen} sind
darunter als Specialfälle enthalten.

Es sind also in dem Gebiete der \gesperrt{complexen Zahlen} $x+yi$ die
Grundrechnungsarten unbegrenzt auszuführen (mit Ausnahme der Division
durch Null), und die Rechnung mit den reellen Zahlen ist ein
Specialfall davon.

Man stellt die complexen Zahlen $z=x+yi$ nach \gesperrt{Gauss}
geometrisch durch die Punkte einer Ebene dar, indem man ein
rechtwinkliges Coordinatensystem zu Grunde legt und den Punkt mit den
Coordinaten $x,y$ als Bild des Zahlwerthes $z$ betrachtet. Die Punkte
der $x$-Axe stellen in der oben besprochenen Weise die reellen Zahlen
$x$ dar. Die Punkte der $y$-Axe sind die Bilder der rein imaginären
Zahlen $yi$. Der Coordinatenanfangspunkt ist das Bild der Zahl $0$.
Der Radiusvector vom Nullpunkte nach dem Punkte $z$ hat den Zahlwerth
$\varrho=\sqrt{x^2+y^2}$, der der \gesperrt{absolute Werth} oder der
Betrag oder, nach älterer Ausdrucksweise, der Modulus der complexen
Zahl $z$ genannt wird.
\endgroup
